\section{Treinamento e Regularização Geométrica}
\label{sec:treinamento}

A perda principal continua sendo a entropia cruzada autoregressiva:
\begin{equation}
  \mathcal{L}_{CE} = -\sum_t \log p(x_{t+1}\mid x_{\leq t}).
\end{equation}
Sobre ela são adicionadas perdas geométricas. A regularização da métrica
penaliza normas excessivas de $U(x)$:
\begin{equation}
  \mathcal{L}_{metric}=\mathbb{E}\|U(x)\|_F^2.
\end{equation}
A diversidade métrica incentiva $U(x)$ a variar ao longo dos tokens, evitando
uma métrica constante. A perda de ortogonalidade aproxima os eixos semânticos
low-rank de direções menos redundantes.

Como os primeiros experimentos mostraram colapso das coordenadas do manifold
próximo de $0{,}5$, foi introduzida uma perda de variância:
\begin{equation}
  \mathcal{L}_{var} =
  \frac{1}{d_m}\sum_{\ell=1}^{d_m}
  \max(0, s_0-\operatorname{std}(x_\ell))^2.
\end{equation}
Essa perda recebe gradiente até o projetor $q\rightarrow q_m$, ao contrário das
perdas aplicadas somente à MetricNet.

Para $d_m=4$, adicionamos uma regularização toroidal explícita. Os pares
$(x_0,x_1)$ e $(x_2,x_3)$ são tratados como dois círculos em torno de $0{,}5$.
Com $z_1=(x_0,x_1)-0{,}5$ e $z_2=(x_2,x_3)-0{,}5$, a perda contém:
\begin{itemize}
  \item termo radial, que aproxima $\|z_1\|$ e $\|z_2\|$ de um raio alvo;
  \item termo de cobertura angular, que força médias vetoriais próximas de
  zero;
  \item termo de isotropia, que evita elipses degeneradas;
  \item termo de independência entre os dois ciclos;
  \item termo de segunda harmônica, que evita cobertura bilobada.
\end{itemize}
O baseline estável usou raio alvo $0{,}35$ e pesos internos ajustados
experimentalmente para equilibrar cobertura, espessura radial e independência.

Para testar se a topologia pode ser sustentada sem imposição permanente, a
implementação também suporta um currículo linear para os pesos geométricos. Em
particular, \texttt{lambda\_torus\_start} e \texttt{lambda\_torus\_end}
permitem iniciar o treino com regularização toroidal forte e reduzi-la até zero
ao longo dos passos. Além disso, foram adicionadas perdas menos prescritivas:
\texttt{intrinsic\_dimension\_loss}, que favorece uma dimensão efetiva alvo
por espectro de covariância, e \texttt{coverage\_entropy\_loss}, que incentiva
cobertura marginal sem especificar ciclos, raios ou pares de coordenadas.
